\documentclass[11pt]{article}

\usepackage[a4paper, margin=1in]{geometry}
\usepackage[T1]{fontenc}
\usepackage[utf8]{inputenc}
\usepackage{amsmath, amssymb}
\usepackage{booktabs}
\usepackage{graphicx}
\usepackage{hyperref}
\usepackage{xcolor}
\hypersetup{colorlinks=true, linkcolor=blue, urlcolor=blue, citecolor=blue}
\usepackage{authblk}
\usepackage{microtype}
\usepackage{enumitem}
\setlist[itemize]{leftmargin=1.4em}
\setlist[enumerate]{leftmargin=1.6em}

\newcommand{\grad}{\nabla}
\newcommand{\half}{\tfrac{1}{2}}

\title{\textbf{Proton--Electron Cosmology}}

\author[1]{Claes Johnson}
\affil[1]{KTH Royal Institute of Technology, Stockholm, Sweden \\ \texttt{claesjohnson@gmail.com}}
\date{\today \\[0.4em]
\small\textit{Mathematical theory and formulation by the author; manuscript developed} \\
\small\textit{in collaboration with Claude (Anthropic). See Section~\ref{sec:collab}.}}

\begin{document}
\maketitle

\begin{abstract}
We sketch an exploratory cosmology built from the RealQM / RealNucleus program, in which all
structure is Coulombic (no strong force, no elementary neutron). The arena is a \textbf{given 3D
Euclidean space} (in the spirit of Many-Minds Relativity) --- no physical medium, only the coordinate
system and the Laplacian. On it live \textbf{two fluctuating potentials}, a gravitational $\varphi_g$ and
an electric $\varphi_E$. Mass and charge are not postulated but read
off as the curvature of the potentials, $\rho=\grad^2\varphi$, which makes both \emph{variable in
sign} and \emph{net-zero} ($\int\grad^2\varphi\,dV=0$): the universe carries no net mass and no net
charge, the latter giving equal numbers of protons and electrons with no baryogenesis. The Laplacian
amplifies small scales ($\rho_k\propto k^2\varphi_k$). A positive-mass region collapses
gravitationally into a hot dense state; its mass splits into \textbf{two positive-mass species ---
protons (heavier) and electrons (lighter), in equal numbers} --- and an electric-potential fluctuation
\emph{in this region} then charges them ($+$ and $-$); deuterons form, and the released binding energy
drives a \textbf{nuclear-powered bounce} into expansion. Electric structure and light exist \emph{only
in this positive-mass region}; the negative-mass regions are chargeless --- literally dark --- and repel
as \textbf{dark energy}. The electron is taken to
have \textbf{two temperature-selected size-phases} --- compact (inside nuclei, ${\sim}$MeV) and
expanded (atomic, ${\sim}$eV) --- so a single cooling curve orders nuclei then atoms, with
neutron $=$ proton $+$ compact electron, the weak interaction as the size-transition, neutron stars
as forced compactification, and the primordial helium fraction set by a transition temperature
$T_\ast$. We state the picture, its one quantitative handle ($T_\ast$), and its load-bearing open
problems --- chiefly \emph{what sets the nuclear scale}, whether the \emph{bounce energy budget}
beats gravity, and whether the $\grad^2\varphi$ \emph{fluctuation spectrum} can match the CMB.

\medskip
\noindent\textbf{Status.} This is an exploratory proposal, not established cosmology. It departs
sharply from the Standard Model and $\Lambda$CDM and must eventually be reconciled with them; the
open problems in Section~\ref{sec:open} are load-bearing. Read it as ``what follows \emph{if}.''

\medskip
\noindent\textbf{Keywords:} cosmology; bounce cosmology; induced gravity; proton--electron model of
the nucleus; dark energy; primordial nucleosynthesis; RealQM; multiphase Coulomb continuum mechanics.
\end{abstract}

\section{Introduction}
The RealQM program formulates quantum mechanics as multiphase 3D continuum mechanics --- unit-charge
densities on non-overlapping domains separated by Bernoulli free boundaries, minimising the Coulomb
energy~\cite{realqm}. Its nuclear extension, RealNucleus, revives the pre-1932 proton--electron
picture (neutron $=$ proton $+$ electron) and reproduces nuclear binding from pure Coulomb
interactions with no strong-force postulate~\cite{realnucleus}. This note asks what cosmology results
if the \emph{same} ingredients --- Coulomb attraction, a multiphase free-boundary vacuum, and
particles built from charge --- are taken as the starting point of the universe. The result is a
collapse-and-bounce cosmology with dark energy and a unified treatment of nuclear and atomic
structure built in. We develop it as a chain of postulates and consequences, and we keep the open
problems explicit.

\section{The arena: a given 3D Euclidean space}\label{sec:substrate}
Following \emph{Many-Minds Relativity}~\cite{mmr}, we take the arena to be a plain \textbf{3D Euclidean
coordinate system} --- absolute space, simply \emph{given}: not Einstein's relative spacetime, and
\emph{not} a physical elastic medium. On it live scalar potentials (a gravitational $\varphi_g$ and an
electric $\varphi_E$), and the only operator needed is the \textbf{Laplacian $\grad^2$}. Mass and charge
are not separate ingredients --- they are the Laplacian (curvature) of the potentials,
\[ \rho_{\text{mass}} = \frac{1}{4\pi G}\,\grad^2\varphi_g, \qquad
   \rho_{\text{charge}} = -\varepsilon_0\,\grad^2\varphi_E . \]
Nothing is presupposed but the coordinate space and the operator --- no pre-stressed medium, no tension to
justify. (This replaces the earlier tensioned-vacuum picture with a smaller, cleaner given.)

\textbf{The Laplacian's large multiplicative effect.} $\grad^2$ multiplies each Fourier mode by $k^2$
($\rho_k\propto k^2\varphi_k$), so a gentle, small-scale ripple in a potential becomes a \emph{large}
density contrast --- a strong multiplicative amplification at short wavelengths (the ``inflation'' of the
density field). A tiny fluctuation of $\varphi$ thus seeds strongly-varying mass and charge on small scales.

\textbf{Mass is not electromagnetic.} The electromagnetic sector carries \emph{charge}; \emph{mass is
inertia} --- the localization, i.e.\ the frequency / inverse size of the matter wave, the dispersion gap
$\omega_0=mc^2/\hbar$ --- a separate, gravitational quantity (signed, hence not the positive-definite EM
field energy). Electromagnetism and gravitation meet \emph{only through energy} ($E=mc^2$), embodied in
matter (a charge welded to an inertia). MMR also gives $E=mc^2$ from a wave momentum $P=mc$
(\cite{mmr}, ch.~16), mass--energy equivalence without Lorentz kinematics; the present construction needs
only the Euclidean arena and the Laplacian, and is otherwise \emph{non-relativistic --- no $c$ enters} the
electrostatic/gravitational structure.

\emph{The irreducible given} is then only a Euclidean coordinate space and the potentials on it --- a far
smaller assumption than a pre-stressed elastic solid. Why there is a space with potentials at all is the
boundary where physics hands off to metaphysics; we state it plainly rather than smuggle it in.

\section{Two fluctuating potentials: mass and charge as curvature}\label{sec:origin}
The fundamental objects are two potentials on the substrate of Section~\ref{sec:substrate} --- a
gravitational $\varphi_g$ and an electric $\varphi_E$ --- each carrying a small-scale fluctuation.
Density is not assumed; it is the Laplacian of the potential (Poisson, run backwards):
\begin{equation}
\rho_{\text{mass}} = \frac{1}{4\pi G}\,\grad^2\varphi_g,
\qquad
\rho_{\text{charge}} = -\varepsilon_0\,\grad^2\varphi_E .
\end{equation}
Three consequences follow at once.
\begin{enumerate}
\item \textbf{Variable sign.} $\rho$ is positive where $\varphi$ is concave and negative where convex.
Gravity thus has $\pm$mass and electromagnetism $\pm$charge, both \emph{derived} from one potential
each, not postulated as separate ingredients.
\item \textbf{Net zero.} For any localised fluctuation, $\int\grad^2\varphi\,dV$ is a surface term
that vanishes, so $\int\rho\,dV=0$. The universe carries \textbf{zero net mass and zero net charge}
automatically. Charge neutrality means equal numbers of protons and electrons --- $\#p=\#e$ --- with
\emph{no baryogenesis required}; mass neutrality means the $\pm$mass regions balance.
\item \textbf{Small-scale amplification.} For a Fourier mode, $\rho_k\propto k^2\varphi_k$: the
Laplacian amplifies short wavelengths by $k^2$, so a gentle potential ripple becomes a large density
contrast. We refer to this as the ``inflation'' of the density field; it is amplification, not the
exponential spatial expansion of standard inflation (Section~\ref{sec:open}).
\end{enumerate}
The Laplacian's $k^2$ amplification first makes many \emph{small} $\pm$mass regions; same-sign mass
attracts and opposite repels, so they \textbf{accumulate into larger positive-mass and negative-mass
regions}. Charge is added \emph{later}, by an electric-potential fluctuation within the positive-mass
region that becomes matter (Section~\ref{sec:bounce}).

\section{The collapse--bounce and dark energy}\label{sec:bounce}
The cosmic sequence runs as a self-regulating collapse--bounce:
\begin{enumerate}
\item A positive-mass region \textbf{collapses gravitationally}, converting potential energy into
heat: the hot dense state is \emph{generated}, not assumed.
\item In it the mass \textbf{splits into two positive-mass species} --- the heavier \textbf{protons},
the lighter \textbf{electrons} (the mass difference set here, gravitationally) --- and a fluctuation of
the electric potential \emph{in this region only} then charges them ($+$ protons, $-$ electrons), in
equal numbers. (This gravitational mass is distinct from the electrostatic kinetic coefficient $\kappa$,
which treats both as charge clouds for atomic/nuclear structure.)
\item \textbf{Deuterons form and release energy}, and that energy drives expansion --- reversing the
collapse. A nuclear-powered bounce.
\item Expansion $\to$ cooling $\to$ atoms. The negative-mass regions are left \emph{chargeless} --- no
electromagnetics, no light, hence literally \textbf{dark} --- acting on the positive-mass region only by
gravity as \textbf{dark energy}; electric structure and light live only in the visible (positive-mass)
region.
\end{enumerate}
Two features stand out: the picture \emph{avoids an initial singularity} (a collapse-and-bounce, not a
Big Bang from nothing, in the spirit of bounce cosmologies~\cite{bounce}), and it \emph{unifies the
expansion drivers} --- the initial expansion from the nuclear bounce, the ongoing acceleration from
negative-mass repulsion. The load-bearing test is the energy budget: the released nuclear binding
energy must exceed the \emph{gravitational} binding of the collapsing region for the bounce to occur
(Section~\ref{sec:open}).

\section{Proton--electron charges and the two electron sizes}\label{sec:scales}
We take protons ($+1$) and electrons ($-1$) to be \textbf{equal in mass} at the level of the charge
fluctuation. In the proton--electron picture a neutron is $p+e$, so any neutral atom $(A,Z)$ contains
$A$ protons and $A$ electrons: $A-Z$ compact inside the nucleus (one per neutron) and $Z$ expanded in
orbitals. Equal numbers of protons and electrons everywhere is consistent with the net-zero charge of
Section~\ref{sec:origin}.

The electron is taken to occupy \textbf{two size-phases}, its extent adapting to the Coulomb well:
\emph{compact} in a deep well (inside a nucleus, ${\sim}$fm, ${\sim}$MeV) and \emph{expanded} in a
shallow one (atomic, ${\sim}$\AA, ${\sim}$eV). The same particle then spans the nuclear and atomic
energy scales, a ratio ${\sim}10^{5}$--$10^{6}$. Chemistry and nuclear physics become one Coulomb
free-boundary mechanism at two scales: a molecule is nuclei glued by \emph{expanded} (valence)
electrons; a nucleus is protons glued by \emph{compact} electrons. What physically sets the
separation is the central open problem (Section~\ref{sec:open}).

\paragraph{Why the gap is $\alpha^2$.} The electron has three characteristic lengths --- the three ways
to build a length from $(\hbar,m,c,e)$, each differing by one factor of $\alpha=e^2/\hbar c$:
\begin{equation}
\underbrace{a_0=\tfrac{\hbar^2}{me^2}\approx0.5\,\text{\AA}}_{\text{atomic }(\hbar,e)}
\;\xrightarrow{\;\times\alpha\;}\;
\underbrace{\bar\lambda_C=\tfrac{\hbar}{mc}\approx386\,\text{fm}}_{\text{Compton }(\hbar,c)}
\;\xrightarrow{\;\times\alpha\;}\;
\underbrace{r_e=\tfrac{e^2}{mc^2}\approx2.8\,\text{fm}}_{\text{classical / nuclear }(e,c)},
\qquad a_0:\bar\lambda_C:r_e = 1:\alpha:\alpha^2 .
\end{equation}
The two electron \emph{phases} are the ends --- expanded (atomic, $a_0$; non-relativistic Coulomb) and
compact (nuclear, $r_e$; electromagnetic mass). The middle length, the \emph{Compton wavelength}
($\hbar,c$; the quantum-relativistic scale where confinement costs ${\sim}mc^2$), is the quantum rung
between them, which is why the nuclear/atomic gap is $\alpha^2$ (two steps of $\alpha$). Dropping
$\hbar$ carries one from the atomic to the classical scale; $\alpha$ is the price of each step. (The
value of $\alpha$ itself is an input, Section~\ref{sec:open}.)

\paragraph{The two sizes, set by charge and proton size --- without mass.} Described by charge and a
kinetic coefficient $\kappa$ alone (the electrostatic structure needs \emph{no mass}), the \emph{same}
electron takes two forms. \textbf{Outside (atomic):} attracted to one proton, its density must taper to
zero in vacuum; a varying profile carries kinetic energy $\kappa/L^2$, forcing a large cloud --- size set
by the proton's \emph{charge}, $L=a_0$. \textbf{Inside (nuclear):} caged by protons with a multiphase
\emph{free boundary} --- its density does \emph{not} drop to zero but hands off to the proton density. A
flat caged density then has $\nabla\psi=0$, hence \emph{zero kinetic energy}: no kinetic penalty for being
small. Its size is set by the proton's \emph{finite size} $R_p$ (the cage floor that stops collapse), its
binding Coulomb at that size ($\sim$MeV at fm). So the crucial factor is \textbf{proton charge vs proton
size}: the ratio is $a_0/R_p\approx1/\alpha^2$ (since $R_p\approx$ the classical radius $\approx$ a few fm).
\textbf{The nuclear scale follows from Coulomb plus the proton's finite size, via a zero-kinetic-energy
flat caged electron --- no relativistic effect required.} (The kinetic obstacle $\kappa/L^2$ applies only
to a \emph{vacuum} electron forced to taper to zero, not to a caged one whose density stays finite at the
free boundary.) Mass plays no part: it is a separate, \emph{gravitational} quantity (the matter-wave's
inverse size / dispersion gap).

\paragraph{Inputs of the electrostatic model.} The whole atom/nucleus structure needs just three:
\textbf{(i)} the charge / Coulomb potential $e^2$, \textbf{(ii)} the electron kinetic coefficient $\kappa$,
and \textbf{(iii)} the proton's finite size $R_p$. From them: the atomic scale $a_0=2\kappa/e^2$ (from
$\kappa,e^2$), and the nuclear scale $R_p$ with binding ${\sim}e^2/R_p$ (from $R_p,e^2$), ratio $a_0/R_p$.
\emph{No mass, no $c$, no strong force, and $\alpha$ is not fundamental} --- the nuclear/atomic ratio is
$a_0/R_p$, which equals $1/\alpha^2$ only because $R_p$ happens to be the classical radius. Mass (inertia)
and gravity form a separate sector, coupling in only through energy ($E=mc^2$).

\paragraph{What $\alpha$ signifies.} The fine-structure constant then emerges as a pure \emph{geometric
ratio} of the two scales,
\[ \alpha^2 = \frac{R_p}{a_0} = \frac{\text{proton (nuclear) size}}{\text{atomic size}}, \qquad
\alpha = \sqrt{R_p/a_0}, \]
--- how much smaller the nucleus is than the atom. \emph{Not a coupling constant, and no $c$ required.} It
equals the textbook $\alpha=e^2/\hbar c$ only because the proton sits at the classical-radius scale
($R_p\approx\alpha^2 a_0\approx$ a few fm). So the model reframes ``why $\alpha\approx1/137$'' as ``why is
the proton ${\sim}$fm while the atom is ${\sim}$\AA{}'' --- relocating the question to the proton size (an
input), and interpreting $\alpha$ itself as the proton-to-atom size ratio.

\paragraph{Empirical check.} Observed sizes: the H atom $\approx a_0=0.53$ \AA{} $\approx 53{,}000$ fm; the
deuteron nucleus $\approx 2.1$ fm (rms). Their ratio $\approx 4\times10^{-5}\approx 1/25{,}000$ matches
$\alpha^2\approx 5\times10^{-5}\approx1/18{,}800$ to within ${\sim}30\%$, with the deuteron (${\sim}2.1$ fm)
sitting essentially at the classical electron radius ($r_e=\alpha^2 a_0=2.8$ fm). So the $\alpha^2$
nuclear/atomic size relation is confirmed to order of magnitude --- the H-atom/deuteron-nucleus size ratio
comes out as $1/\alpha^2$, from charge, $\hbar$ and $c$ alone, with no strong force. (The deuteron is
unusually loosely bound, so this is the scale/ratio being right, not a precise fit.)

\section{Temperature, the cosmic sequence, and the weak interaction}\label{sec:temp}
A bound state survives once the bath can no longer ionise it ($kT<$ binding). The two electron phases
therefore switch on at two temperatures set by their two binding energies, forcing the order
plasma $\to$ nuclei $\to$ atoms (Table~\ref{tab:epochs}). The temperature ratio between the two epochs
(${\sim}10^6$) equals the energy-scale ratio --- one fact, not two.

\begin{table}[h]
\centering
\begin{tabular}{lll}
\toprule
epoch & process & temperature \\
\midrule
plasma & free $\pm$ charges & $>10^{10}$ K \\
nuclear condensation & compact $e$ captured $\to$ nuclei & ${\sim}10^{10}$ K ($kT{\sim}$MeV) \\
recombination & expanded $e$ captured $\to$ atoms; CMB & ${\sim}3000$ K ($kT{\sim}$eV) \\
\bottomrule
\end{tabular}
\caption{The two condensation epochs, separated by ${\sim}10^6$ in temperature.}
\label{tab:epochs}
\end{table}

The compact electron is precisely the electron inside a neutron: $\text{neutron}=p+\text{compact }e$.
The two weak processes are then the two directions of the size-transition, driven by
temperature/density: \emph{compactify} ($p+e\to n$, electron capture) and \emph{expand}
($n\to p+e$, $\beta$-decay). This lands on real astrophysics: crushing matter forces electrons
compact $\to$ wholesale $p\to n \to$ a neutron star (neutronisation in core collapse). The cosmic
neutron/proton ratio at freeze-out is the compact/free electron ratio at $T_\ast$.

\section{The quantitative handle: $T_\ast$ and the helium fraction}\label{sec:helium}
The postulate predicts a transition temperature $T_\ast$ at the compactification energy. The fraction
of electrons that compactify by freeze-out fixes the cosmic neutron/proton ratio and hence the
\textbf{primordial helium fraction}; matching the observed ${\sim}25\%$ He ($n/p\approx1/7$) is the
decisive quantitative test. In standard physics freeze-out is at $kT\approx0.8$ MeV
($T\approx10^{10}$ K), tied to the $n$--$p$ mass difference $1.29$ MeV, so the prediction is
$T_\ast\sim$ MeV $\sim10^{10}$ K. Equivalently, the universe's small:big (compact:expanded) electron
ratio equals the $N{:}Z$ of primordial matter, ${\approx}1{:}7$, dominated by hydrogen (which has no
compact electron).

A neutron is a proton with a compact electron, so the neutron/proton ratio is the compact/free
electron ratio at freeze-out, $(n/p)_f=e^{-E_\ast/kT_f}$, with the compactification gap $E_\ast$
playing the role of the $n$--$p$ mass difference; free neutrons partly expand back ($\beta$-decay,
factor $d$) before nucleosynthesis, and essentially all survivors end in He-4:
\begin{equation}
(n/p)_{\text{nuc}}=\frac{(n/p)_f\,d}{1+(n/p)_f\,(1-d)},
\qquad
Y=\frac{2\,(n/p)_{\text{nuc}}}{1+(n/p)_{\text{nuc}}} .
\end{equation}
With the model's single MeV-scale handle ($E_\ast\approx1.29$ MeV, $T_f\approx0.8$ MeV so
$T_\ast\sim10^{10}$ K, $d\approx0.74$): $(n/p)_f=0.199\to(n/p)_{\text{nuc}}=0.140\to\boxed{Y=0.245}$,
the observed primordial helium. A single compactification energy thus reproduces the ${\sim}25\%$
helium --- the model's \emph{first falsifiable quantitative success}. (Caveat: $T_f$ is, in standard
BBN, derived from the weak-rate-vs-expansion balance; here it is a second input, so the genuine
prediction is the ratio $E_\ast/kT_f\approx1.6$ with $E_\ast\sim1.3$ MeV.)

\section{Energetics: where the binding energy goes}\label{sec:energetics}
Forming a deuteron releases its binding (${\sim}$MeV). Because the binding is Coulombic --- an
electron dropping from the expanded into the compact phase --- the energy is released as photons,
exactly as an atomic electron radiates when it falls to a lower level, scaled up ${\sim}10^5$ to
${\sim}$MeV gammas. Equivalently it appears as the deuteron's mass defect ($E=mc^2$): the deuteron is
lighter than its constituents by binding$/c^2$. Cosmologically it is dumped into a photon bath with
${\sim}10^9$ photons per baryon, a ${\sim}10^{-8}$ perturbation that redshifts to insignificance ---
the universe is not noticeably reheated.

The process is \textbf{not self-sustaining}: the released gammas feed the same bath that re-dissociates
fresh deuterons while $kT\gtrsim$ binding, so net formation proceeds only as the universe \emph{cools}
below $T_\ast$ (the deuterium bottleneck). This is the opposite of barrier-limited stellar fusion
(favoured by heating, energy release sustaining the high $T$ it needs); deuteron \emph{formation} is
exothermic capture (favoured by cooling) and self-quenches at high $T$. Proton density sets the
\emph{yields} (as the baryon density does in standard nucleosynthesis), not ignition. The one
self-sustaining channel is \emph{gravitational}: in a collapsing dense region, rising density forces
electron capture in a runaway sustained by gravity --- neutron-star neutronisation.

\paragraph{Radiation (deterministic, no photons).} The light of this cosmology --- the thermal radiation of
the hot dense state and its relic, the microwave background --- is handled within RealQM \emph{deterministically},
without photons or statistics, by the \emph{interaction matter--radiation} developed in~\cite{realqm}. Matter
and radiation share one wave-equation structure: each electron density and the radiation field obey a wave
equation augmented by an Abraham--Lorentz radiation-reaction term (energy radiated by accelerating charge) and a
viscous term (coherent energy converted to heat at small scales). This yields the Rayleigh--Jeans law at low
frequency and a temperature-dependent high-frequency cutoff --- Wien's displacement law, and the Planck spectrum
--- with the ultraviolet catastrophe removed by a finite resolution scale rather than by quantisation, and the
second law of radiation (heat flows hot $\to$ cold) following directly from the PDE. So the cosmic thermal
radiation belongs to the same deterministic continuum program, not a separate quantised sector.

\section{Open problems}\label{sec:open}
\begin{enumerate}
\item \textbf{The nuclear scale.} What sets the ${\sim}10^4$--$10^5$ nuclear/atomic separation?
\emph{Candidate (preferred):} take the \textbf{proton's finite size $R_p$ (${\sim}$fm) as the input
nuclear length}. An electron caged by protons with a multiphase \emph{free boundary} (density finite, not
dropping to zero) can be \emph{flat} $\Rightarrow$ \emph{zero kinetic energy}, so it sits at the cage
(proton) size with no kinetic penalty; binding is Coulomb (${\sim}$MeV at fm). The size ratio to the
atomic (charge-set) electron is $a_0/R_p\approx1/\alpha^2$ --- no relativity needed. \emph{Alternative:}
the relativistic \emph{classical radius} $r=\alpha^2 a_0\approx2.8$ fm (Coulomb self-energy $=$ rest
energy, in the sense of Many-Minds Relativity~\cite{mmr}), giving the same ratio. Open which primitive
(proton size vs.\ relativistic), and whether the solver realises the flat zero-KE cage --- direct runs so
far were inconclusive (non-convergence at small cage size). The kinetic obstacle $\kappa/L^2$ applies only
to a \emph{vacuum} electron, not a caged one.
\item \textbf{The bounce energy budget.} Does the released nuclear binding energy exceed the
gravitational binding of the collapsing region, so the bounce of Section~\ref{sec:bounce} actually
reverses collapse? At large scales gravity typically dominates nuclear energy (why supernova cores
collapse rather than explode on fusion alone). The inequality $E_{\text{nuclear}}\gtrsim
E_{\text{grav}}$ at the relevant scale must be checked.
\item \textbf{The fluctuation spectrum and $\varphi$ dynamics.} $\rho=\grad^2\varphi$ gives a blue
($k^2$-amplified) density spectrum, whereas the CMB and large-scale structure fit a near
scale-invariant one ($n_s\approx0.96$); the potential must carry a compensating red spectrum.
Moreover $\rho=\grad^2\varphi$ is a definition, not yet a law of motion --- $\varphi$ needs its own
field equation.
\item \textbf{The mass ratio $m_p/m_e=1836$.} The origin's mass-split makes the proton heavier than the
electron (gravitationally), so the mass difference is now \emph{built in} --- but its \emph{value}
($1836$) is not derived (nor is it in the Standard Model: proton mass ${\sim}99\%$ QCD binding, electron
mass from the Higgs). A subtlety to reconcile: the quantitative nuclear binding (He-4/deuteron $13.1$
vs.\ $12.7$) treats protons as \emph{charge clouds} with an electron-like kinetic coefficient, so the
\emph{gravitational} mass (heavier proton) is distinct from the \emph{electrostatic} kinetic coefficient
$\kappa$ (both charge clouds) --- the relation between the two is open. Numerological aside:
$m_p/m_e\approx6\pi^5=1836.1$.
\item \textbf{The neutrino.} $\beta$-decay has a continuous spectrum --- the 1930 objection that sank
the proton--electron neutron and forced Pauli's neutrino~\cite{pauli}. The ``electron expands'' step
needs a home for that energy and spin.
\item \textbf{Precision data and the rest.} BBN abundances and the CMB acoustic peaks are
percent-level successes the sequence must reproduce; antimatter and proton substructure must be
accounted for in a two-ingredient ontology.
\end{enumerate}

\section{Computational status}\label{sec:comp}
What the RealQM/RealNucleus engine has shown, versus what remains. \textbf{Done:} $p$-$e$-$p$ binds as
a stable configuration; the He-4/deuteron binding \emph{ratio} comes out $13.1$ vs.\ experiment
$12.7$~\cite{realnucleus}; the lepton-mass scaling law $R\propto1/m_{\text{glue}}$ is reproduced, with
muon-catalysed fusion (the muonic molecular ion $dd\mu$ at ${\sim}512$ fm) as the observed proof that a
heavy negative glue binds two protons; and the helium-fraction calculation of
Section~\ref{sec:helium} gives $Y=0.245$ from a single MeV-scale handle, matching the observed
primordial helium. \textbf{Open:} the two-scale test (does one engine give MeV \emph{and} eV in the
same units? equal-mass Coulomb returns ${\sim}1$, confirming a scale ingredient is needed); and a
cooling-quench ``mini Big Bang'' using the engine's temperature and Brownian-dynamics machinery to
enact the two condensation epochs.

\section{Complexity from the proton--electron size asymmetry}\label{sec:complexity}
The world is built not on a \emph{charge} imbalance (charge is exactly balanced, $\pm e$) but on a
\emph{size} imbalance: the proton has a \textbf{fixed, small size} (a rigid anchor) while the electron's
size is \textbf{variable} --- it adapts to its well, spreading to ${\sim}$\AA{} in an atom and compacting
to ${\sim}$fm in a cage. The fixed proton supplies the anchors; the variable electron \emph{spans scales}
--- spreading to bind atoms and molecules, compacting to bind nuclei --- producing the nested multi-scale
hierarchy that \emph{is} complexity. With a fixed size for both, only one scale exists and nothing complex
is built (both small $\to$ only nuclei; both large $\to$ only atoms). And $\alpha$ measures exactly this
asymmetry: $\alpha^2=R_p/a_0$ is the \emph{inverse dynamic range} of the electron --- a factor
$1/\alpha^2\approx18{,}800$ between its nuclear and atomic phases --- so \textbf{the smallness of $\alpha$
is the engine of complexity}. Equal fixed sizes would build nothing.

\section{Collaboration}\label{sec:collab}
The mathematical theory and the cosmological proposal are the author's. The manuscript was developed
in dialogue with Claude (Anthropic), which also assisted in articulating the argument, surfacing the
quantitative constraints and open problems, and preparing the supporting interactive simulations of
the RealQM/RealNucleus program. We present the human--AI collaboration transparently as part of the
working method.

\paragraph{One-sentence summary.} On a given 3D Euclidean space, if mass is built first from a
gravitational-potential fluctuation (protons heavier, electrons lighter) and charge is then added by an
electric-potential fluctuation in the positive-mass region --- the negative-mass regions left chargeless
and dark --- and the electron has two temperature-selected sizes, then nuclei, atoms, the weak
interaction, neutron stars, dark energy, and the primordial helium fraction follow from one cooling curve
--- a unifying picture whose survival
hinges on a single hard question: what sets the nuclear scale, if not the strong force.

\paragraph{Three questions to close.} Why bring in the strong force at all, when RealNucleus binds nuclei
by Coulomb alone (neutron $=$ p$+$e), using and needing no strong interaction? How crucial is the
\emph{exact} value of the nuclear/atomic scale difference $\alpha^2\approx1/18{,}800$? And would the world
look essentially the same with a scale difference ten times smaller --- nuclei ten times larger relative to
atoms --- still hierarchical, still complex? The picture suggests that what matters for a structured
universe is that a \emph{size asymmetry exists at all} (a rigid small proton, a fluid variable electron),
far more than its precise magnitude --- so the world may be robust, not finely tuned, in the value of
$\alpha^2$.

\begin{thebibliography}{9}
\bibitem{realqm} C.\ Johnson, \emph{RealQM: Quantum Mechanics as Multiphase 3D Continuum Mechanics}
(arXiv article; see its \emph{Interaction matter--radiation} section for the deterministic matter--light /
blackbody model) and interactive gallery, \url{https://claes542.github.io/RealMolecule/}.
\bibitem{realnucleus} C.\ Johnson, \emph{RealNucleus: nuclear binding from a pure-Coulomb
proton--electron model}, manuscript, same repository.
\bibitem{mmr} C.\ Johnson, \emph{Many-Minds Relativity} --- special relativity as medium/aether
physics: $c$ as the wave speed of a real light-carrying medium, mass as field self-energy, $P=mc
\Rightarrow E=mc^2$ (ch.~16), Galilean kinematics, and a critique of Einstein's special relativity.
\bibitem{sakharov} A.\ D.\ Sakharov, \emph{Vacuum quantum fluctuations in curved space and the theory
of gravitation}, Dokl.\ Akad.\ Nauk SSSR \textbf{177} (1967) 70 [Sov.\ Phys.\ Dokl.\ \textbf{12},
1040].
\bibitem{bounce} M.\ Novello and S.\ E.\ P.\ Bergliaffa, \emph{Bouncing Cosmologies}, Phys.\ Rep.\
\textbf{463} (2008) 127.
\bibitem{pauli} W.\ Pauli, open letter to the Tübingen group (1930), proposing the neutrino to save
energy conservation in $\beta$-decay.
\bibitem{kolbturner} E.\ W.\ Kolb and M.\ S.\ Turner, \emph{The Early Universe}, Addison-Wesley
(1990).
\end{thebibliography}

\end{document}
